\section{PART I --- LINEAR ALGEBRA (THEORETICAL)}

\subsection{1. Vector Spaces}
\begin{itemize}
\item Fields and definitions; vector spaces and subspaces
\item Span, linear independence; bases and dimension
\item Sum and direct sum of subspaces
\item Quotient spaces; first isomorphism theorem for vector spaces
\item Why this matters: all of PDE, signal processing, and ML live in vector spaces
\end{itemize}

\subsection{2. Linear Maps}
\begin{itemize}
\item Maps as the primary object, matrices as secondary
\item Null space and range; rank-nullity theorem
\item Isomorphisms; space of linear maps $L(V,W)$
\item Quotient constructions; exact sequences (preview)
\item Applications: linear systems, PDE operators, neural network layers
\end{itemize}

\subsection{3. Matrices as Representations}
\begin{itemize}
\item Matrix of a linear map in a basis; change of basis; similarity
\item Row reduction as a computational tool; rank from the map perspective
\item Equivalence vs.\ similarity vs.\ congruence
\item Connection to numerical linear algebra (Part II)
\end{itemize}

\subsection{4. Dual Spaces}
\begin{itemize}
\item Linear functionals; dual space $V^*$; examples
\item Dual basis; double dual and canonical isomorphism
\item Transpose of a linear map; annihilators
\item Reflexivity; connection to functional analysis and weak formulations in PDE
\end{itemize}

\subsection{5. Eigentheory}
\begin{itemize}
\item Eigenvalues and eigenvectors; characteristic polynomial
\item Cayley-Hamilton theorem; consequences for matrix polynomials
\item Invariant subspaces; triangularization over $\mathbb{C}$
\item Generalized eigenvectors and Jordan normal form
\item Minimal polynomial; relation to Jordan structure
\item Applications: stability of ODEs, Google PageRank, PCA, Markov chains
\end{itemize}

\subsection{6. Inner Product Spaces}
\begin{itemize}
\item Inner products and norms; Cauchy-Schwarz
\item Gram-Schmidt orthogonalization
\item Orthogonal complements and projections; best approximation
\item Adjoint $T^*$; normal, self-adjoint, unitary, and orthogonal operators
\item Applications: least squares, QR factorization, signal decomposition
\end{itemize}

\subsection{7. Spectral Theory}
\begin{itemize}
\item Spectral theorem: real symmetric (orthogonal diagonalization)
\item Spectral theorem: normal operators over $\mathbb{C}$
\item Positive semidefinite operators; Cholesky as spectral factorization
\item Singular value decomposition: existence, uniqueness, geometry
\item Applications: PCA, dimensionality reduction, pseudoinverse, graph analysis
\end{itemize}

\subsection{8. Multilinear Algebra}
\begin{itemize}
\item Bilinear forms and quadratic forms; Sylvester's law of inertia
\item Tensor products of vector spaces; universal property
\item Exterior algebra; determinant as top exterior power; Pfaffian
\item Symmetric powers; applications in physics, differential geometry, ML
\item Tensor products as the foundation for all tensor decompositions (Part II)
\end{itemize}

\subsection{9. Modules and Canonical Forms}
\begin{itemize}
\item Modules over a ring; free modules; connection to vector spaces
\item Structure theorem for finitely generated modules over a PID
\item Jordan and rational canonical forms as corollaries
\item Why this matters: unifies eigentheory; also underpins homological methods
\end{itemize}

\section{PART II --- LINEAR ALGEBRA (COMPUTATIONAL)}

\subsection{10. Floating Point and Numerical Stability}
\begin{itemize}
\item IEEE arithmetic; machine epsilon; rounding models
\item Condition number: of a matrix and of a problem
\item Forward and backward error analysis; Wilkinson's approach
\item Catastrophic cancellation; why naive algorithms fail
\end{itemize}

\subsection{11. LU and Cholesky Factorization}
\begin{itemize}
\item Gaussian elimination as LU factorization
\item Partial and complete pivoting; stability analysis
\item Cholesky for symmetric positive definite systems
\item Solving $Ax = b$ stably; cost analysis
\item Applications: scientific simulation, Bayesian linear systems
\end{itemize}

\subsection{12. QR Factorization}
\begin{itemize}
\item Modified Gram-Schmidt; why classical fails numerically
\item Householder reflections; Givens rotations
\item QR for least squares; normal equations vs.\ QR
\item Rank-deficient least squares; pseudoinverse; regularization
\item Applications: regression, overdetermined systems, data fitting
\end{itemize}

\subsection{13. Singular Value Decomposition (Computational)}
\begin{itemize}
\item Golub-Reinsch algorithm; bidiagonalization
\item Truncated SVD; optimal low-rank approximation (Eckart-Young)
\item Applications: PCA, image compression, latent semantic analysis, pseudoinverse
\item Connection to randomized methods (Section 18)
\end{itemize}

\subsection{14. Tensor Decompositions and Tensor Networks}
\begin{itemize}
\item CP (CANDECOMP/PARAFAC) decomposition; alternating least squares
\item Tucker decomposition; multilinear SVD (HOSVD)
\item Tensor train (TT) / matrix product states; bond dimension
\item NP-hardness of tensor rank; uniqueness conditions
\item Applications: neural network compression, quantum many-body, high-dimensional PDEs
\end{itemize}

\subsection{15. Eigenvalue Algorithms}
\begin{itemize}
\item Power iteration; inverse iteration; shift-and-invert
\item QR algorithm with shifts and deflation
\item Symmetric QR with Wilkinson shift; divide-and-conquer
\item Krylov subspace methods: Lanczos (symmetric), Arnoldi (general)
\item Applications: Google PageRank, structural mechanics, Schr\"odinger operators
\end{itemize}

\subsection{16. Iterative Methods for Linear Systems}
\begin{itemize}
\item Krylov methods: CG (SPD), MINRES (symmetric), GMRES (general), BiCGSTAB
\item Convergence theory; connection to polynomial approximation of eigenvalue distribution
\item Preconditioning: ILU, incomplete Cholesky, algebraic multigrid (conceptual)
\item CG as optimization; connection to Course 5
\end{itemize}

\subsection{17. Sparse Matrices}
\begin{itemize}
\item Storage formats: CSR, CSC, COO; arithmetic on sparse formats
\item Sparse direct methods: fill-in, reordering (AMD, nested dissection)
\item Sparse iterative methods in practice; choosing a solver
\item Applications: large-scale PDE, graph problems, network analysis
\end{itemize}

\subsection{18. Randomized Linear Algebra}
\begin{itemize}
\item Randomized SVD (Halko-Martinsson-Tropp); single-pass variant
\item Sketching and random projections; Johnson-Lindenstrauss lemma
\item Randomized least squares; sketch-and-solve; sketch preconditioners
\item Kaczmarz algorithm as SGD on normal equations
\item Applications: large-scale ML, streaming data, massive linear systems
\end{itemize}

\subsection{19. Spectral Graph Theory}
\begin{itemize}
\item Graph Laplacian (combinatorial and normalized); spectrum examples
\item Cheeger inequality: algebraic connectivity vs.\ graph expansion
\item Random walks on graphs; stationary distributions; mixing times
\item Spectral clustering; manifold learning (Laplacian eigenmaps, diffusion maps)
\item PageRank; graph neural networks as spectral operators
\end{itemize}

\subsection{20. Structured Matrices and Fast Algorithms}
\begin{itemize}
\item Toeplitz, circulant, Hankel matrices; displacement rank
\item FFT as diagonalization of circulant matrices
\item Hierarchical matrices (H-matrices): structure and arithmetic (brief)
\item Fast multipole method (conceptual); applications to $N$-body and integral equations
\end{itemize}

\subsection{21. Matrix Functions and Matrix Calculus}
\begin{itemize}
\item Matrix exponential: Pad\'e approximation and scaling-squaring
\item Other functions: square root, logarithm, sign; Krylov approximation
\item Fr\'echet derivative of matrix functions
\item Matrix calculus: gradients and Hessians of scalar matrix functions
\item Applications: ODE solvers, network analysis, optimization on matrix manifolds
\end{itemize}

\section{PART III --- ABSTRACT ALGEBRA IN SERVICE OF APPLICATIONS}

\subsection{22. Groups and Symmetry}
\begin{itemize}
\item Definitions and examples: symmetric, cyclic, dihedral, matrix groups
\item Subgroups; Lagrange's theorem; cosets; normal subgroups
\item Why this matters: symmetry is a modeling primitive in physics, chemistry, and ML
\end{itemize}

\subsection{23. Group Homomorphisms and Quotients}
\begin{itemize}
\item Homomorphisms and isomorphisms
\item Quotient groups; isomorphism theorems
\item Direct products; semidirect products
\item Applications: symmetry breaking, gauge theory, signal processing on groups
\end{itemize}

\subsection{24. Group Actions and Symmetry in Computation}
\begin{itemize}
\item Actions, orbits, stabilizers; orbit-stabilizer theorem
\item Burnside's lemma; counting with symmetry
\item Cayley's theorem; conjugacy classes; class equation
\item Applications: molecular symmetry, enumeration, equivariant architectures
\end{itemize}

\subsection{25. Sylow Theory and Group Structure}
\begin{itemize}
\item Sylow theorems: existence, conjugacy, number (statements and proofs)
\item Classification of groups of small orders; direct applications
\item Simple groups: $A_5$ is simple (one proof)
\item Classification of finite abelian groups
\item Applied payoff: understanding symmetry groups of physical systems; underpins representation theory
\end{itemize}

\subsection{26. Rings and Arithmetic}
\begin{itemize}
\item Definitions; examples ($\mathbb{Z}$, polynomial rings, matrix rings, $\mathbb{Z}/n\mathbb{Z}$)
\item Ideals, quotient rings, ring homomorphisms; isomorphism theorems
\item Product rings and Chinese Remainder Theorem (with algorithm)
\item Prime and maximal ideals
\item Applications: modular arithmetic, error-correcting codes, cryptography
\end{itemize}

\subsection{27. Divisibility, Factorization, and Polynomial Rings}
\begin{itemize}
\item ED $\to$ PID $\to$ UFD: the hierarchy; emphasis on payoffs
\item Gaussian integers; applications to number-theoretic algorithms
\item Polynomial rings: Eisenstein, reduction mod $p$
\item Applications: irreducibility testing, polynomial system solving
\end{itemize}

\subsection{28. Field Theory and Finite Fields}
\begin{itemize}
\item Field extensions; degree $[K:F]$; algebraic elements and minimal polynomial
\item Splitting fields; normal and separable extensions
\item Finite fields $\mathrm{GF}(p^n)$: structure, existence, primitive element
\item Applications: coding theory (Reed-Solomon, LDPC), cryptography (AES, elliptic curves), pseudorandom generation
\end{itemize}

\subsection{29. Galois Theory}
\begin{itemize}
\item Galois group; fundamental theorem of Galois theory
\item Computing Galois groups (examples over $\mathbb{Q}$)
\item Solvable groups; solvability by radicals; quintic insolvability
\item Cyclotomic fields; roots of unity
\item Applied payoff: underpins algebraic coding theory; explains why generic degree-5 polynomial systems require numerical rather than closed-form solution
\end{itemize}

\subsection{30. Modules}
\begin{itemize}
\item Modules over a ring; submodules; quotients; homomorphisms
\item Free modules; exact sequences; Snake lemma
\item Structure theorem for finitely generated modules over a PID
\item Tensor product of modules; Hom functor; adjunction
\item Applications: unifies linear algebra over arbitrary rings; foundation for homological algebra
\end{itemize}

\subsection{31. Representation Theory of Finite Groups}
\begin{itemize}
\item Group algebras $k[G]$; representations as $G$-modules
\item Maschke's theorem: complete reducibility in characteristic 0
\item Irreducible representations; Schur's lemma
\item Characters and character tables; orthogonality relations
\item Applications: Burnside's theorem ($p^a q^b$ solvable), Fourier analysis on finite groups, design of symmetric algorithms
\end{itemize}

\subsection{32. Harmonic Analysis on Groups}
\begin{itemize}
\item Pontryagin duality for locally compact abelian groups
\item Peter-Weyl theorem for compact groups
\item Spherical harmonics as irreducible representations of $SO(3)$
\item Generalized Fourier analysis; convolution on groups
\item Applications: fast algorithms on groups, non-commutative Fourier analysis, physics
\end{itemize}

\subsection{33. Equivariant Maps and Geometric Deep Learning}
\begin{itemize}
\item Equivariant and invariant linear maps; Schur's lemma computationally
\item $SO(3)$ and $SE(3)$ representations; spherical harmonics explicitly
\item Equivariant neural network architectures (E3NN, NequIP)
\item Gauge equivariance in physics-inspired models
\item Applications: molecular property prediction, protein structure, robotics
\end{itemize}

\subsection{34. Gr\"obner Bases and Polynomial Systems}
\begin{itemize}
\item Polynomial ideals; term orderings; Buchberger's algorithm
\item Reduced Gr\"obner bases; ideal membership; elimination theory
\item Solving polynomial systems; robotics kinematics; computer vision
\item Software: SymPy, Macaulay2, Singular
\end{itemize}

\subsection{35. Commutative Algebra and Algebraic Geometry (Brief)}
\begin{itemize}
\item Noetherian rings; Hilbert basis theorem
\item Localization; local rings; geometric meaning
\item Prime spectrum $\mathrm{Spec}(R)$; Hilbert's Nullstellensatz
\item Applied payoff: foundations of algebraic statistics, sum-of-squares optimization, polynomial optimization
\end{itemize}

\subsection{36. Homological Algebra and Topological Data Analysis}
\begin{itemize}
\item Chain complexes and homology; long exact sequences
\item Functors Hom and $\otimes$; derived functors Ext and Tor
\item Simplicial complexes as concrete homological objects
\item Persistent homology: filtrations, barcodes, persistence diagrams, stability theorem
\item Software: Ripser, Gudhi, Giotto-TDA
\item Applications: shape analysis, materials, neuroscience, time series topology
\end{itemize}
